\documentclass[]{../../../../tex/classes/styledArticle}
\usepackage{../../../../tex/packages/hebrewSupport}

\author{שחר פרץ}
\title{\textit{היסטוריה 39}}
\begin{document}
	\maketitle
	\section{תנועת המרי}
	תנועת המרי היא ההתארגנות הראשונה של הישוב היהודי, כישוב, במאבק נגד הבריטי. הוקם באוקטובר 1945 ופעלה במשך 10 חודשים. תנועת המרי היא ''ארגון גג`` שיש בו אצ''ל, לח''י, הגנה ופלמ''ח. בראש תנועת המרי הייתה ועדה בשם ועדת X. כל ארגון נשאר עצמאי, עם נשקו ועם האידיאולוגיה שלו. מטרת הועדה היה להעביר מידע ולתאם פעולות. יש שלושה דרכי פעולה: העפלה, התיישבות, ופעולות צבאיות. 
	
	תנועת המרי פעלות צבאיות. אף אחד לא היה צריך את ועדת X כדי לתאם העפלה. דרך הפעולה והתפיסה היא פגיעה בבריטים במתקנים צבאיים ואסטרטגיים. אפילו האגף הדרומי במלון המלך דוד היה מפקדה בריטי, ולא מלון. 
	
	הרעיון – כוחנו באחדותנו. לפני התנועה, המחתרות נאבקו בינהן. היו תקופות (נקראו ''סזון``, בצרפתית עונת הציד) בה ההגנה והפלמ''ח הסגירו פעילים של האצ''ל והלח''י. 
	
	בהקמת התנועה, בן גוריון ישר את ההגנה והפלמ''ח עם תפיסה יותר קיצונית ממנו. אחת הסיבות היא האכזבה מממשלת הללייבור, ונאום בוין. בין היתר הייתה כוונה להשפיע על תוצאות הודעה האנגלו־אמריקאית. 
	
	הפעולות של התנועה הן סימליות, ונעדו לפגוע בחלקים הלא־צודקים לדעת היישוב של השלטון הבריטי בארץ. לדוגמה: 
	\begin{itemize}
		\item שחרור מחנה המעפילים בעתלית באוקטובר 1945. 200 מעפילים שוחררו ופוזרו ברחבי הארץ מבלי להתפס מחדש. זו הפעולה היחידה שבה הבריטים לא מצליחים לתפוס את המעפילים. זו אינה פעולת העפלה, אלא פעולה צבאית לכל דבר. 
		\item ''ליל הרכבות`` – גם הוא באוקטובר, בו חובלו 153 נקודות ברשת מסילות הברזל ברחבי הארץ. הופעלו בו כל יחידות הפלמ''ח. זו הפעולה היחידה שמבוצעת על־ידי כל המחתרות ביחד. 
		\item הטבעת אוניות משמר – פלוגות מיוחדות של הפלמ''ח הטביעו שלוש אוניות משמר בריטיות. 
		\item ליל הגשרים – מבצע מרכולת (17.6.2025) שם המפלמ''ח סגר את כל הגשרים המובילים לארץ ישראל תוך פיצוצם בלילה אחד. המבצע עבר כמעט כולו בהצלחה (10 גשרים פוצצו ללא בעיה), פרט בגשר הרכב (גשר הזיו) נתגלו החבלנים, ו־14 מהם נהרגו והגשר לא פוצץ. 
	\end{itemize}
	
	אומנם היה תיאום בין התנועות, אבל לרוב שלושת המחתרות פעלו בנפרד. רק בליל הרכבות לא היה זה המצב. 
	
	בתקופה הזו הבריטים עשו פעולות ומעצרים בנסיון למנוע את זה. אחרי ליל הגשרים, ''הגיעו מים עד נפש``, והם פתחו במבצע שנקרא ''מבצע השבת השחורה``. הם סגרו את הגבולות, ניתקו קווי תקשורת וביצעו עוצר בערים המרכזיות. הם חיפשו ועצרו לוחמים, בינהם מפקדים בכירים כמו רבין (נעצרו בערך 2000 לוחמים), וכן פוליטיקאים. הבריטים גם החרימו נשק ביישוב. בקיבוצים היה המון נשק. אחד הקיבוצים האלו היה יגור (עד היום מדי פעם יש כתבות על בתים ישנים שנהרסים, ומוצאים בהם סליק מתקופת המנדט). ביגור היה סליק ענק, שהבריטים גילו כמעט בקושי. הם דפקו על האדמה ומצאו כך את הסליק. החרימו ביגור ''המון נשק``. 
	
	בערך חודש אחרי השבת השחורה האצ''ל מבצע את הפיצוץ במלון המלך דוד. בעבר, זה היה ברור לכלם, שהאצ''ל עשה את הפעולה על דעת עצמו ולא קיבל אישור על ועדת X. במהלך השנים התברר שזה לא היה בדיוק כך – ועדת X כנראה אישרה את הפעולה באופן עקרוני, אך רק בשעות הלילה בעת שלא היו כבר אנשים (הרי רוצים לפוצץ את המפקדה, ולא עובדים בה בלילה). האצ''ל אכן נתן התראה והבריטים סירבו להתפנות. בסופו של דבר נהרגו 91 אנשים. 
	
	ע''פ חוק עזר עירוני בירושליים, הבתים אמורים להיבנות באבנים ירושלמיות. היו לוקחים בטון ופשוט מצפים אותו. בעבר זה אשכרה היה בנוי באבנים ירושלמיות, והקירות במלון המלך דוד היו ברוחב ענק מאבנים. כדי לפוצץ את זה היה צריך כמויות ביזאריות של חומר נפץ, והכמות הגדולה של חומר הנפץ גרמה לקריסה של האגף, ולהרוגים רבים שחלקם אפילו סתם הלכו ברחוב ליד ולא היו קשורים. 
	
	שני האירועים האלו גורמים לפירוק המועצמה: בשבת השחורה, הנהגת היישוב מבינה שזו לא הדרך. הבריטים יותר חזקים ויש להם יותר כלים. 
	
	כאשר פוצץ מלון המלך דוד, נהרגו אנשים. זה סוקר רע בעולם ופגע בדעת הקהל. הפגיעה בתדמית של המאבק ביישוב מאוד משמעותית. זו הסיבה העיקרית לפירוק – ההגנה והפלמ''ח הם המנהיגים של היישוב והתנועה הציונית, ולא יכולים לקחת חלק בזה. 
	
	המחלוקת האידיאולוגית הייתה כל הזמן, ולא נעלמה בעת קיום תנועת המרי. עם זאת, הם הסכימו להיות תחת המטריה הזו בגלל האינטרסים המשותפים. 
	
	(הבהרה: כל אותה העת יש מאבק צמוד של ההגנה+פלמ''ח וכן אצ''ל+לחי. תנועת המרי אפשרה שיתוף פעולה בין המאבקים הצמודים האלו, אבל המאבק האידיאולוגי נמשך גם בעת קיומה). 
	
	
	
	
	\ndoc
\end{document}